一个 64 位哈希，碰撞概率有多小？

常见的回答是``小到可以忽略''。但只要记录条数超过 \(2^{64}\)，碰撞就不是小概率事件，而是必然事件——抽屉只有这么多，东西比抽屉多，必然有两件挤在一起。这个论证不需要任何概率假设，也不依赖数据分布，它只用到一件事：可能取值的个数是有限的。

同一个有限性还有另一面。三十个城市的巡回路线也是有限多条，可它有 \(29!/2\) 种，把宇宙年龄拿来枚举都不够。有限意味着``理论上数得清''，完全不意味着``实际上算得完''。

离散数学就在这两面之间工作。它一边利用有限性榨取必然结论——鸽巢原理、握手引理、Lagrange 定理、泵引理都是同一种手法：数量对不上，所以某件事非发生不可；另一边又要正面对付有限带来的爆炸，而唯一的出路是\textbf{在对象里找出结构}，让搜索不必逐个进行。

\begin{center}
\begin{tikzpicture}[x=1cm,y=1cm,font=\small,>=stealth]
  \draw[black!45,fill=green!8] (0,1.1) rectangle (6.2,3.3);
  \node[align=center] at (3.1,2.75) {\textbf{有限 \(\Rightarrow\) 必然}};
  \node[align=center,font=\footnotesize] at (3.1,2.0) {鸽巢原理、握手引理\\Lagrange 定理、泵引理};
  \node[align=center,font=\footnotesize,black!65] at (3.1,1.4) {数量对不上，某件事非发生不可};
  \draw[black!45,fill=red!7] (6.9,1.1) rectangle (13.1,3.3);
  \node[align=center] at (10,2.75) {\textbf{有限 \(\Rightarrow\) 爆炸}};
  \node[align=center,font=\footnotesize] at (10,2.0) {\(n!\) 条路线、\(2^n\) 个子集\\哈密顿回路、可满足性};
  \node[align=center,font=\footnotesize,black!65] at (10,1.4) {数得清，却算不完};
  \draw[black!45,fill=blue!8] (3.2,-0.9) rectangle (9.9,0.35);
  \node[align=center] at (6.55,0.0) {两边的出路是同一个：\textbf{在对象里找出结构}};
  \node[align=center,font=\footnotesize,black!65] at (6.55,-0.62) {双射、递推、图、同余、多项式、群、格};
  \draw[->,black!55] (3.1,1.05) -- (5.0,0.4);
  \draw[->,black!55] (10,1.05) -- (8.1,0.4);
\end{tikzpicture}

\emph{九章各自处理一类结构，但都在回答同一个问题：这些对象是怎样生成的？看清生成方式，计数、判定和构造才有可能绕开逐个枚举。}
\end{center}

\subsection{五条贯穿始终的判断}

\textbf{数一个集合，就是把它对应到已经会数的集合上。}对应可以是一一的（双射），可以是整齐的多对一（除法原理），也可以换个方向重新清点同一批东西（双重计数）。排列、组合、重复排列不是三条需要分别记住的公式，它们是同一个除法原理在不同对象上的实例。对应建立不起来时才退到鸽巢（只断言碰撞会发生）或容斥（把重叠算干净）。

\textbf{表面相似的问题，难度可以相差极远。}走遍每条边（欧拉回路）有一个线性时间的度数判据；走遍每个顶点（哈密顿回路）是 NP 完全的。两句话的表述几乎对称，可判定性却天差地别。差别的位置很具体：欧拉回路不存在时你能当场给出简短的否定证书（某个顶点度数为奇），哈密顿回路不存在时没人知道怎样给出这种证书。立项之前先判断需求落在哪一边，比动手实现更要紧。

\textbf{代数结构按``保留哪些运算规律''分层，每加一条要求，换来一种能力。}群只管一种可逆运算；环放进两种，但不再保证除法，于是模 6 下会出现 \(2\times3\equiv0\) 这样的零因子；域补上非零元可逆，插值与解方程才成立。格走的是另一条路线，把偏序里的上确界与下确界当成两种运算。读这些结构时不要背公理，要问每条公理挡住了什么失效。

\textbf{对称可以被代数化，然后反过来约束设计。}群是``哪些改动不改变本质''的写法；等变要求写成方程 \(W\rho_X(g)=\rho_Y(g)W\) 之后，解它就能得到具体的层结构——平移群解出循环矩阵（卷积），置换群解出 \(aI+b\mathbf{1}\mathbf{1}^{\trans}\)（集合上的聚合）。这不是事后类比，是参数共享的代数来源。

\textbf{有限的记忆决定能力的边界。}有限自动机必须把无限多的输入前缀归并成有限个等价类，因此记不住任意长的计数——泵引理把这句话变成了可用的否定工具。同一种论证方式在别处反复出现：抽屉有限、状态有限、位宽有限，能装下的区分度就有上限。

\subsection{九章怎么安排}

主线是三章：\hyperref[ch:027]{``组合计数：先描述对象怎样生成，再决定怎样数''}给出计数的基本动作，\hyperref[ch:029]{``图论：局部邻接怎样形成全局网络''}给出关系建模的通用语言，\hyperref[ch:031]{``初等数论：整除结构怎样变成可计算算法''}给出有限算术的地基。这三章后面被引用得最多。

\hyperref[ch:028]{``递推关系：从局部生成规则恢复整体公式''}紧接组合计数，处理``规模 \(n\) 的答案由更小规模组成''这类问题，是分治与动态规划的共同语言；它与线性代数的特征值是同一件事的两种写法。

\hyperref[ch:031]{``编码与多项式方法：用结构化冗余对抗丢失和错误''}把数论与多项式汇合起来，并通向信息论。\hyperref[ch:032]{``群与对称：可逆操作怎样组织不变性''}是几何深度学习的数学地基，其中``不变与等变''一篇值得单独读。

\hyperref[ch:033]{``布尔代数与自动机：有限规则怎样处理逻辑与序列''}篇幅短，读它的理由是那条能力边界，不是化简技巧。\hyperref[ch:034]{``环与格：两种运算怎样组织代数与次序''}偏抽象，但格在类型系统、静态分析和 CRDT 里有直接对应，可按需回补。\hyperref[ch:035]{``群表示论：对称的线性结构''}需要线性代数在手，第一次读离散数学时不必把它当作终点。

阅读前最好熟悉\hyperref[ch:004]{``数学语言：读清一句话究竟断言了什么''}里的量词、集合与归纳法。读完之后，组合计数直接进入\hyperref[ch:036]{``概率论：在不确定性中进行可检验的推理''}——概率的第一步就是把样本空间的结构数清楚。

拿到一个新问题时，这一部分希望你先问三句：对象是什么、约束是什么、这个结构逼出了什么必然结论。把这三句问顺了，选哪种工具是随后自然的事；反过来，先记住 \(\binom nk\) 和一堆算法名字，遇到没见过的问题仍然无从下手。
